\chapter{Szintillationsdetektor}
\section{Theoretische Grundlagen}
\subsection{Radioaktiver Zerfall und $\gamma$-Strahlung}
Instabile Atomkerne wandeln sich durch radioaktiven Zerfall in energetisch günstigere Zustände um. Bei vielen Zerfallsprozessen verbleibt der Tochterkern zunächst in einem angeregten Zustand, bevor er unter Emission eines Gamma-Photons in einen niedrigeren Energiezustand übergeht. Diese Photonen weisen diskrete Energien auf, die für das jeweilige Nuklid charakteristisch sind.\\

Zu den im Experiment verwendeten wichtigen Quellen gehören Cs-137, Co-60 und Eu-152, deren Zerfallsarten wohldefinierte Gamma-Linien liefern.
\subsection{Wechselwirkung von $\gamma$-Strahlung mit Materie}

\subsection{NaI(Tl)-Szintillationsdetektor}
\subsection{Signale und Spektren}
\subsection{Linienbreite und Energieauflösung}
\subsection{Detektoreffizienz}
\subsection{Energiekalibrierung}
